\section{有理数域}

\begin{definition}{有理数域}{}
	我们称$\Z$的分式域为有理数域,记作$\Q$,其元素称为有理数.
\end{definition}

\begin{definition}{正有理数集}{}
	定义$\Q_{\geq 0}=\left\{\dfrac{a}{b}\in\Q\middle|a\in\N,b\in\N_{>0}\right\},\Q_{\leq 0}=-\Q_{\geq 0},\Q_{>0}=\Q_{\geq 0}\setminus\left\{0\right\},\Q_{<0}=-\Q_{>0}$.
\end{definition}

\begin{proposition}{}{}
	$\Q_{\geq 0}$是$\Q$的子环,$\Q=\Q_{<0}\sqcup\left\{0\right\}\sqcup\Q_{>0}$,在同构的意义下$\Q_{\geq 0}\cap\Z=\N$.
\end{proposition}

\begin{definition}{$\Q$上的全序}{}
	定义$\Q$上的关系$\leq$如下:$\forall x,y\in\Q\left(x\leq y\iff y-x\in\Q_{\geq 0}\right)$.
\end{definition}

\begin{proposition}{$\Q$上的标准全序的性质}{}
	\begin{enumerate}
		\item $\leq$是$\Q$上的全序,配备该全序后$\Q$和$\Z$是序同构的.
		\item (加法相容性)$\forall x,x^{\prime},y,y^{\prime}\in\Q\left(x\leq y\wedge x^{\prime}\wedge y^{\prime}\implies x+x^{\prime}\leq y+y^{\prime}\right)$.
		\item (加法相容性)$\forall x,,y\in\Q\left(x\leq y\implies \forall z\in\Q\left(\left(z\in\Q_{\geq 0}\implies xz\leq yz\right)\right)\wedge\left(z\in\Q_{\leq 0}\implies xz\geq yz\right)\right)$.
	\end{enumerate}
\end{proposition}

\begin{proposition}{$\Q^{\times}$的分解}{}
	我们有典范的群同构$\Q^{\times}\simeq\left\{-1,1\right\}\times\Q_{>0}$.
\end{proposition}

\begin{definition}{$\Q$上的绝对值}{}
	记$\varphi:\Q^{\times}\to\Q_{>0}\times\left\{-1,1\right\}$为典范同构.
	\begin{enumerate}
		\item 定义$\left|\cdot\right|:\Q\to\Q_{>0}$如下:对$x\in\Q^{\times}$,$\left|x\right|$是$\varphi\left(x\right)$在$\Q_{>0}$中的分量;对$x=0$,$\left|x\right|=0$.
		\item 定义$\sgn:\Q\to\left\{-1,0,1\right\},x\mapsto
		\begin{cases}
			1,&x\in\Q_{>0},\\
			0,&x=0,\\
			-1,&x\in\Q_{<0}.
		\end{cases}$
	\end{enumerate}
	我们称$\left|\cdot\right|$是$\Q$上的绝对值,称$\sgn$为$\Q$的符号函数.
\end{definition}


